\section{Dependence of the gate law at width 1024}\label{sec:dep}

Section~5 of the companion composes a Gaussian observation tier with a conditional lower tier of gap $\delta$. Its \S2.2
explains that spectral independence would supply $\delta$, provided influence is controlled under every pinning. The
unpinned influence radius $\eta=\lambda_{\max}(\corr g)-1$ is a necessary condition for that, and it is measurable
(Table~\ref{tab:gap}).

\begin{itemize}[leftmargin=1.5em]
\item $\eta$ is about $2$ at layer $1$, matching the Marchenko--Pastur prediction $3\cdot2/\pi=1.91$. It rises to $6$--$8$
in the deep layers. The gate correlation spectrum is broad, with participation ratio $520$--$720$ out of $1024$. Even
unpinned, the law is far from the $\eta<1$ regime.
\item Conditioning on $k$ principal components of the pre-activations lowers $\eta$ only slowly: at $k=16$ it is still
$4$--$5$ at depth, and at $k=128$ it is $1.2$--$1.6$. The Marchenko--Pastur noise edge at these sample sizes is
$0.18$--$0.27$.
\item Gates do not decorrelate across depth. $\|\corr(g_\ell,g_{\ell+d})\|_{\rm op}$ is $2.4$ at layer $1$ and $6$--$7$ at
depth, with essentially no decay over $d=1,2,4$.
\item The closure's level-1 depth gain $\prod_{k>\ell}g_k$, with $g_k=2\,\overline{\Phi(\alpha_k)^2}$, is $0.03$ from layer
$1$, $0.31$ from layer $8$ and $0.6$--$0.9$ over the last four layers. Unresolved late-layer structure is not forgotten
before the output.
\end{itemize}

The wall-chain bound of Proposition~\ref{prop:lin} is reported per layer from the measured gate covariance and wall rates.
It is $0.20$--$0.33$ per unit rotation angle on every layer, against $2/\pi=0.64$ for independent gates. Removing $16$
principal components raises it to $0.23$--$0.65$; removing $128$ raises it to $0.33$--$1.7$. Gates with large $|\alpha|$
have both small variance and slow walls, which offsets the larger $\eta$ at depth. The physical single-layer chain is
therefore never degenerate on these networks, while the information heat-bath walk has no comparable guarantee.

The level process relaxes at rate $0.65$ at layer $1$, the independent-gate value $2/\pi$, and at $2$--$3$ at depth. That
stiffness is what pins $r$ at $n/2$. Whether the joint chain over all $nL$ gates, the companion's full-history graph,
keeps this rate is a separate question, because gate dependence does not decay across depth. Section~\ref{sec:full}
measures it.

\section{The full-history wall graph}\label{sec:full}

The companion's geometric graph uses complete gate histories. Proposition~\ref{prop:lin} applies to any set of gates, so
it bounds that graph by the rate-normalised covariance of all $nL$ gates at once (\texttt{scripts/s24\_fullgates.py}).
On net 0 there are $14{,}748$ live gates in $N=2^{16}$ inputs.

\begin{itemize}[leftmargin=1.5em]
\item The stacked gate correlation has many large eigenvalues: $71,63,59,56,55,53$ against a noise edge of $2.2$. The
rate-normalised covariance has top eigenvalues $41.6,37.8,35.9,34.1,32.9,32.4$. Hence
\[
\gap(\text{full-history wall chain})\le 1/41.6=0.024\ \text{per unit rotation angle},
\]
about nine times below any single layer.
\item The bound falls steadily with the number of layers taken together. The best windows give $0.21$ for one layer,
$0.12$--$0.17$ for two, $0.064$--$0.089$ for four, $0.038$--$0.043$ for eight and $0.024$ for all sixteen.
\item The slowest mode is spread over the whole depth. Each layer holds $3.5$--$8.7\%$ of the top eigenvector, with most in
layers $5$--$12$.
\end{itemize}

Cross-depth coherence is what makes the local full-history chain slow, not any single layer. In the companion's terms
this is the quantitative content of ``local adjacency does not inherit the dense gap'' for the real networks. The actual
rotation process decorrelates within an angle of $\pi/2$, while its local Markov shadow needs about $1/0.024\approx40$.
A two-tier construction must therefore place these coherent cross-depth modes in its global observation tier, and the
count of such modes sets the dimension of that tier. The count of such modes was not measured; that run was stopped in favour of theory (Stage 25).

\section{The terminal ellipsoid at width 1024}\label{sec:ell}

For hidden states equal to inputs, the companion's operator $\mathcal B$ in eq.~(44) has the nonzero spectrum of $S/n$,
where $S=\cov F(X)$ is the output covariance. A first-pass error $e=\mu-b$ can be resolved in the eigenbasis of $S$
(\texttt{scripts/s24\_ellipsoid.py}; net 0 also from the larger run behind \texttt{scripts/s24\_terminal.py}).
Table~\ref{tab:ell} gives the results.

\begin{itemize}[leftmargin=1.5em]
\item The visible spectrum is concentrated. The participation ratio of $S$ is about $51$. The top $16$, $64$ and $256$
eigenvectors carry $42\%$, $75\%$ and $95\%$ of the trace, and the bottom half carries $1\%$.
\item First-pass errors live where the variance lives. For the closure, $69\%$ of $|e|^2$ is in the top $16$ eigenvectors,
which is the collective tier. CC1 removes most of that, and its remaining error is still $86\%$ inside the top $256$. The
v56 error is the most spread, with $6\%$ in the bottom half. The Rayleigh ratio $e\T Se/(|e|^2\,\overline\lambda)$ is
$66$, $8.3$ and $5.1$ for closure, CC1 and v56 on net 0.
\item Sampling at the floor budget ($M\approx4500$ samples after a CC1-sized pass) resolves this residual only with
information about its direction. We give three oracle bounds for CC1 on net 0, each assuming $e$ is known.

\emph{Orthogonal projections.} The best $\mu=b+P(\bar Y-b)$ has rank at most one, because $ee\T-S/M$ has at most one
positive eigenvalue. Its risk is $1.3\times10^{-7}$, and the projection onto $e/|e|$ gives $1.4\times10^{-7}$. Both are
about four times the floor target $3.1\times10^{-8}$.

\emph{Per-direction choice.} Keeping or dropping each eigendirection with oracle knowledge gains $2\%$ over CC1.

\emph{Oblique filter.} The best linear filter estimates one coefficient along $e$ through the whitened statistic
$e\T S^{+}(\bar Y-b)$. Its risk is $|e|^2/(1+M\,e\T S^{+}e)/n=8.0\times10^{-9}$, below the target, because
$M\,e\T S^{+}e=233$. That signal-to-noise comes from the error's small components along the \emph{quiet} directions,
pooled coherently.

Hence a response direction predicted accurately where $S$ is small would make one Mahalanobis-weighted calibration
sufficient. None of the weight-derived families tried (coherent functions of $a$, propagated coherent functions, tangent
singular subspaces, readout shapes, the CC1$-$closure direction) captures more than $86\%$ of $|e|^2$ even orthogonally.
They are built in the loud directions.

Variance reduction does not change this. Exact-mean input controls of degree $\le k$ reduce variance by at most $1.23$,
$1.52$, $2.10$ and $3.38$ for $k=1,2,4,8$. These are the chaos shares of the depth-16 arc-cosine kernel
(\texttt{theory24/chaos\_spectrum.py}); the measured linear control gives $1.23$. No sampling density helps by more than
$(\E|z|)^2/\E|z|^2\approx1\%$ for the full output, because $|z|$ concentrates.
\end{itemize}

The residual that matters in (44) therefore sits mostly in the same $\sim50$--$250$ modes that carry the output variance,
with a small but statistically decisive part in the quiet directions. The companion's deterministic route has no sampling
variance: one occupation row and operator actions replace Monte Carlo. Its open obligation is a representation of the
operator action that is accurate on the loud modes, with the quiet-direction part as a natural target for a final
whitened calibration.

\section{What the theory and the measurements now say to build}\label{sec:build}

\begin{enumerate}[leftmargin=1.5em]
\item \textbf{The operator.} Use the physical wall chain, or a block or collective refinement of it, not single-gate
heat-bath resampling. Its gap is governed by the rate-normalised gate covariance (Proposition~\ref{prop:lin},
Remark~\ref{rem:near}), and it is bounded away from zero on every single layer of these networks: the level-1 values are
$0.20$--$0.33$ per unit rotation angle. The full-history graph is the open case (Section~\ref{sec:full}).
\item \textbf{The slow block.} About $100$--$250$ modes are needed, by three independent measures. These are the visible
spectrum of $\mathcal B$ (95\% of trace in $256$ modes), the location of CC1's error (86\% in $256$ modes), and the
conditioning that brings the gate law near $\eta\approx1$ ($128$ modes). Those modes are wide but low-rank compared
with $n$.
\item \textbf{The information level.} Record deviations from the closure's reference pattern. The deviation set has
$0.5n$ elements at layer $1$ but only $0.085n$ at layer $16$, and the closure predicts its size. Active cardinality is
pinned at $n/2\pm8$ and carries almost no information.
\item \textbf{Depth.} Gate dependence persists across at least four layers, and the depth gain of the last four layers is
$0.6$--$0.9$. The occupation recurrence cannot truncate late-layer return paths. In the companion's language, $m$ must
grow where the hidden operator's gap is small; that happens in the last layers, not the first.
\item \textbf{What the first pass gives.} The closure is the mean-field wall sum (Proposition~\ref{prop:closure}), and
its wall measures are accurate in aggregate (Rice--Gauss to $0.1$--$0.6\%$). Its error is in the downstream-jump channel,
the conditional expectation of the gated product at a wall. That is the quantity a represented operator must improve.
\end{enumerate}

\appendix
\section{Independent verification of the companion's finite claims}
\texttt{theory24/check\_companion.py}; exact rationals where the claim is finite.

\begin{center}\small
\begin{tabular}{@{}lll@{}}\toprule
claim & check & result\\\midrule
eq.~(7) two-row heat bath & exact spectrum, $\varepsilon=\frac1{10},\frac14,\frac12$ & pass\\
eq.~(13) five-sector example & jumps $(-1,\sqrt2,0,0,1)$, $\mu=1/(2\sqrt\pi)$ to $10^{-10}$; Monte Carlo & pass\\
eq.~(11) $\sqrt{t/\pi}$ constant & half-plane, $t=10^{-2},2.5\cdot10^{-3}$: ratio $0.998$ & pass\\
eq.~(16) finite-noise readout & 3-layer net, $\rho=0.6$: $4.6\times10^{-3}$ relative & pass\\
eq.~(17) one-ReLU moment & $\rho=0.5,0.9$ & pass\\
eq.~(22) posterior gap & three random $M$, 8 cells: gaps $0.83,0.65,0.42\ge0.54,0.26,0.18$ & pass\\
eq.~(28) cycle gap & $m=6,12,40$ & pass\\
Thms 7.1--7.2, eq.~(48), table & exact rationals, $m=0..7$; $\frac12$ without hidden reward, $\frac32$ without clock & pass\\
Thms 8.1--8.2, (39), (43), (44) & 200 random reversible chains & pass\\
eqs.~(3)--(4) & commutator; $D_{r,s}(c)$ counts & pass\\\bottomrule
\end{tabular}
\end{center}

\section{Tables}
\begin{table}[ht]\centering\footnotesize
\caption{Levels and walls per layer, width $1024$ ($N=2^{16}$ inputs and rotation partners per net). Net 0 values; net 0 only (the multi-net run was stopped).}\label{tab:levels}
\resizebox{\textwidth}{!}{\begin{tabular}{@{}rccccccc@{}}\toprule
$\ell$ & $r/n$ & sd$(r)/n$ & excitation$/n$ & $R^2$ next level: signs / values & wall rate & rate sd/mean & Rice--Gauss / measured\\\midrule
1 & 0.500 & 0.0154 & 0.500 [0.500--0.500] & 0.34 / 0.60 & 0.317 [0.317--0.318] & 0.07 & 1.003 [1.000--1.003]\\
2 & 0.504 & 0.0143 & 0.314 [0.300--0.317] & 0.31 / 0.55 & 0.322 [0.312--0.323] & 0.23 & 0.999 [0.999--1.001]\\
3 & 0.503 & 0.0127 & 0.252 [0.248--0.257] & 0.32 / 0.56 & 0.319 [0.314--0.323] & 0.38 & 1.002 [1.001--1.002]\\
4 & 0.497 & 0.0123 & 0.225 [0.221--0.225] & 0.36 / 0.58 & 0.325 [0.320--0.325] & 0.49 & 1.001 [0.997--1.001]\\
5 & 0.498 & 0.0118 & 0.191 [0.185--0.200] & 0.32 / 0.53 & 0.317 [0.316--0.330] & 0.61 & 1.005 [0.996--1.005]\\
6 & 0.502 & 0.0109 & 0.176 [0.173--0.176] & 0.31 / 0.51 & 0.327 [0.320--0.327] & 0.67 & 1.001 [1.000--1.003]\\
7 & 0.507 & 0.0104 & 0.168 [0.147--0.168] & 0.34 / 0.53 & 0.338 [0.301--0.338] & 0.73 & 1.002 [0.997--1.005]\\
8 & 0.464 & 0.0102 & 0.157 [0.144--0.157] & 0.29 / 0.48 & 0.345 [0.312--0.345] & 0.79 & 0.999 [0.996--1.004]\\
9 & 0.497 & 0.0092 & 0.133 [0.126--0.134] & 0.28 / 0.48 & 0.329 [0.299--0.329] & 0.89 & 1.001 [0.999--1.002]\\
10 & 0.528 & 0.0090 & 0.129 [0.117--0.129] & 0.32 / 0.52 & 0.338 [0.296--0.338] & 0.92 & 1.003 [1.001--1.003]\\
11 & 0.518 & 0.0094 & 0.127 [0.103--0.127] & 0.36 / 0.54 & 0.355 [0.277--0.355] & 0.93 & 1.003 [0.997--1.003]\\
12 & 0.493 & 0.0090 & 0.116 [0.103--0.116] & 0.30 / 0.50 & 0.339 [0.288--0.339] & 1.03 & 1.007 [1.002--1.007]\\
13 & 0.516 & 0.0083 & 0.105 [0.086--0.105] & 0.38 / 0.58 & 0.331 [0.256--0.331] & 1.11 & 1.002 [1.002--1.005]\\
14 & 0.522 & 0.0087 & 0.094 [0.086--0.096] & 0.26 / 0.44 & 0.312 [0.269--0.312] & 1.22 & 0.999 [0.999--1.007]\\
15 & 0.500 & 0.0075 & 0.096 [0.088--0.096] & 0.33 / 0.52 & 0.329 [0.284--0.329] & 1.20 & 1.004 [1.004--1.004]\\
16 & 0.511 & 0.0076 & 0.085 [0.085--0.090] & -- & 0.309 [0.297--0.309] & 1.31 & 1.006 [1.006--1.008]\\
\bottomrule\end{tabular}}\end{table}
\begin{table}[ht]\centering\footnotesize
\caption{Dependence of the gate law, single-layer wall-chain bounds (per unit rotation angle) and depth gains, net 0.}\label{tab:gap}
\resizebox{\textwidth}{!}{\begin{tabular}{@{}rcccccccc@{}}\toprule
$\ell$ & $\eta$ & $\eta$, 16 / 128 PCs out & PR$(\corr g)$ & $\lambda_{\rm lin}$ & $\lambda_{\rm lin}$, 16 / 128 PCs out & level relax. & $\|\corr(g_\ell,g_{\ell+1})\|$ & $\prod_{k>\ell}g_k$\\\midrule
1 & 1.96 [1.9--2.0] & 1.77 / 1.53 & 720 & 0.213 [0.210--0.214] & 0.23 / 0.33 & 0.66 & 2.36 & 0.030\\
2 & 2.65 [2.6--2.7] & 2.31 / 1.62 & 692 & 0.216 [0.213--0.219] & 0.25 / 0.41 & 0.80 & 2.97 & 0.050\\
3 & 3.65 [3.6--3.8] & 3.02 / 1.66 & 650 & 0.201 [0.198--0.205] & 0.25 / 0.51 & 0.93 & 3.85 & 0.075\\
4 & 4.57 [4.5--4.7] & 3.53 / 1.60 & 622 & 0.192 [0.187--0.195] & 0.25 / 0.63 & 1.07 & 4.62 & 0.109\\
5 & 5.14 [4.9--5.5] & 3.85 / 1.44 & 612 & 0.197 [0.186--0.202] & 0.27 / 0.60 & 1.21 & 5.29 & 0.145\\
6 & 5.98 [5.7--6.2] & 4.08 / 1.35 & 600 & 0.190 [0.185--0.200] & 0.29 / 0.68 & 1.31 & 5.77 & 0.191\\
7 & 6.01 [5.9--6.2] & 4.32 / 1.31 & 596 & 0.206 [0.204--0.208] & 0.31 / 0.72 & 1.48 & 6.03 & 0.240\\
8 & 6.79 [6.5--7.0] & 4.27 / 1.32 & 577 & 0.202 [0.198--0.205] & 0.34 / 0.63 & 1.60 & 6.54 & 0.309\\
9 & 6.63 [6.4--6.8] & 4.22 / 1.37 & 582 & 0.224 [0.219--0.230] & 0.37 / 0.82 & 1.78 & 6.69 & 0.372\\
10 & 7.03 [6.7--7.3] & 4.27 / 1.37 & 567 & 0.227 [0.219--0.232] & 0.40 / 0.96 & 1.83 & 7.16 & 0.446\\
11 & 7.84 [7.4--8.3] & 4.25 / 1.41 & 545 & 0.219 [0.207--0.226] & 0.43 / 0.81 & 1.88 & 7.65 & 0.525\\
12 & 7.76 [7.4--8.2] & 4.23 / 1.32 & 534 & 0.232 [0.218--0.241] & 0.46 / 0.89 & 1.99 & 7.28 & 0.607\\
13 & 7.13 [6.9--7.5] & 3.97 / 1.28 & 520 & 0.265 [0.249--0.278] & 0.52 / 1.16 & 2.18 & 6.89 & 0.685\\
14 & 7.17 [6.6--8.0] & 4.03 / 1.27 & 517 & 0.277 [0.246--0.315] & 0.54 / 1.15 & 2.31 & 7.15 & 0.780\\
15 & 7.54 [7.1--8.0] & 4.08 / 1.25 & 492 & 0.276 [0.261--0.284] & 0.56 / 1.02 & 2.53 & 7.09 & 0.896\\
16 & 7.16 [6.8--7.7] & 3.90 / 1.25 & 471 & 0.305 [0.273--0.334] & 0.61 / 1.13 & 2.53 & -- & 1.000\\
\bottomrule\end{tabular}}\end{table}
\begin{table}[ht]\centering\footnotesize
\caption{The terminal ellipsoid: visible spectrum of $S=\cov F$ and the share of first-pass errors in the eigenvector bands $[0,16),[16,64),[64,256),[256,512),[512,1024)$; R = Rayleigh ratio $e\T Se/(|e|^2\overline\lambda)$.}\label{tab:ell}
\resizebox{\textwidth}{!}{\begin{tabular}{@{}rcccll@{}}\toprule
net & PR & trace, top 16/64/256 & bottom half & closure: bands; R & CC1: bands; R\\\midrule
0 & 48.9 & 0.418/0.756/0.949 & 0.0099 & 0.689 0.149 0.109 0.037 0.015; 68.4 & 0.209 0.372 0.284 0.104 0.032; 8.3\\
1 & 46.2 & 0.436/0.766/0.950 & 0.0102 & 0.743 0.085 0.120 0.040 0.013; 72.9 & 0.211 0.306 0.337 0.111 0.036; 8.7\\
2 & 50.0 & 0.440/0.772/0.950 & 0.0104 & 0.789 0.082 0.082 0.034 0.012; 70.1 & 0.166 0.356 0.312 0.127 0.039; 8.2\\
3 & 51.5 & 0.425/0.755/0.946 & 0.0117 & 0.671 0.123 0.121 0.065 0.021; 58.7 & 0.248 0.305 0.272 0.132 0.043; 9.7\\
4 & 63.7 & 0.397/0.735/0.937 & 0.0148 & 0.622 0.152 0.150 0.054 0.021; 41.1 & 0.278 0.333 0.254 0.099 0.035; 9.9\\
5 & 44.4 & 0.450/0.773/0.953 & 0.0083 & 0.733 0.147 0.075 0.035 0.011; 74.5 & 0.341 0.310 0.209 0.112 0.027; 11.7\\
6 & 52.5 & 0.422/0.757/0.946 & 0.0109 & 0.619 0.132 0.169 0.060 0.019; 51.7 & 0.318 0.243 0.300 0.108 0.031; 17.6\\
7 & 47.6 & 0.436/0.767/0.953 & 0.0077 & 0.771 0.096 0.093 0.032 0.009; 77.6 & 0.127 0.328 0.387 0.128 0.030; 6.6\\
\bottomrule\end{tabular}}\end{table}

